\begin{abstract}
Deep-learning studies on gallbladder ultrasound usually give accuracy and a handful of heatmaps as proof of interpretability, and rarely test whether either one survives closer scrutiny. This paper does that scrutiny directly, on a ResNet-50 trained for nine-class gallbladder ultrasound classification on the public UIdataGB dataset, and the scrutiny itself is what changes the headline number. An initial split, de-duplicated only by single-orientation image hashing, gave 93.05\% accuracy; a closer audit found that split still let near-identical frames from the same imaging session, and rotation-augmented duplicates the dataset's own release pipeline introduces, cross the train/validation boundary. Correcting for both---clustering images at the session level and by an 8-orientation perceptual hash, then holding out a genuine test set touched only once---drops mean accuracy across four independently seeded models to 72.56\% $\pm$ 2.35\%. We report this corrected figure as the paper's real result, not the inflated one. The main XAI evaluation---five attribution techniques (Grad-CAM, Grad-CAM++, Saliency Maps, Eigen-CAM, and Score-CAM) compared for consistency across replicate models, faithfulness, stability, a parameter-randomization sanity test, a shortcut audit, and calibration---was run on the uncorrected, higher-accuracy split, and Grad-CAM++ and Score-CAM are the most reproducible across seeds there; on faithfulness all three CAM methods perform similarly, and no single method wins on every metric. Because the reproducibility replicates in that battery were trained on disjoint data subsets rather than identical data with only the seed varied, we separately reran the consistency check under a corrected same-data design on the leakage-fixed split, using Grad-CAM as a representative method, together with a null-control test---cosine similarity between untrained models and between synthetic random heatmaps---that the original metric had not been checked against. The corrected same-data consistency score clears both null floors by roughly 2.5$\times$ rather than sitting close to them, so the consistency finding survives this check, though it has not yet been re-verified for all five methods under the corrected design. Grad-CAM++'s output depends on the model's learned weights rather than the architecture alone, and while we find no strong sign the model relies on image-border artefacts specifically, a separate, targeted audit finds a real batch-correlated shortcut: a classifier trained only to identify which acquisition batch a Carcinoma image came from reaches 82.88\% accuracy (8.3$\times$ chance) even after correcting for cluster-level leakage, and a linear probe on the disease classifier's own features recovers batch identity at 74.75\% (4.5$\times$ chance), directly implicating the classifier this paper's explanations describe. The results show these attributions are measurably faithful to the classifier's behaviour; that is not the same as proof the classifier's decisions are clinically reliable, and the batch-shortcut finding above is a concrete reason for caution beyond the usual internal-validation caveats. This remains an internal-validation study: it does not establish multicentre generalizability, patient-level validation, or clinician-confirmed localization, and several methodological gaps identified during external review---including bootstrap intervals for the XAI metrics (now added for the consistency, faithfulness, and stability tables), a full re-run of the five-method battery on the corrected split, and validation of the session/batch inference used in the split correction---remain open and are reported honestly rather than elided.
\end{abstract}
